% Part: many-valued-logic
% Chapter: three-valued-logics
% Section: lukasiewicz

\documentclass[../../../include/open-logic-section]{subfiles}

\begin{document}

\olfileid{mvl}{thr}{luk}

\olsection{منطق~\L ukasiewicz}

من أقدم ما نشر مفصلًا في المنطق المتعدد القيم اقتراح الفيلسوف
البولندي يان~\L ukasiewicz سنة~1921. وكان باعثه مسألة أرسطو في
المعركة البحرية: فقول القائل «ستقع معركة بحرية غدًا» لا يظهر،
\emph{اليوم}، صادقًا ولا كاذبًا؛ لأن ما يعيّن صدقه لم يقع بعد.
فأدخل~\L ukasiewicz لمثل هذه القضايا «المستقبلية العرضية» قيمة ثالثة.
\begin{quote}
يجوز لي، من غير تناقض، أن أفترض أن كوني في وارسو في ساعة معينة من
السنة المقبلة---كظهر يوم~21 ديسمبر---ليس الآن متعينًا بالإيجاب ولا
بالسلب. فيكون وجودي هناك في تلك الساعة ممكنًا، لا ضروريًا. وعلى هذا
الفرض لا يصدق الآن القول «سأكون في وارسو ظهر يوم~21 ديسمبر من السنة
المقبلة»، ولا يكذب. فلو صدق الآن لزم أن يكون وجودي المستقبلي في وارسو
ضروريًا، وذلك خلاف الفرض؛ ولو كذب لزم أن يكون مستحيلًا، وذلك خلافه
أيضًا. فليست القضية الآن صادقة ولا كاذبة، ولا بد لها من قيمة ثالثة
غير «$0$»، أي الكذب، وغير «$1$»، أي الصدق. ولنرمز إليها بـ
«$\frac{1}{2}$»؛ فهي تمثل «الممكن»، وتنضم قيمة ثالثة إلى الصادق
والكاذب.
\end{quote}
وسنرمز هنا إلى قيمة~\L ukasiewicz الثالثة بـ~$\Undef$.
\footnote{مراد~\L ukasiewicz بـ«الممكن» في هذا الموضع غير المعنى
المألوف اليوم؛ فهو الممكن الذي ليس بضروري.}

وعلى هذا الفهم يسهل وضع دوال الصدق لـ~$\lnot$ و$\land$ و$\lor$.
فنفي القضية المستقبلية العرضية مستقبلية عرضية كذلك، ومن ثم
$\tf{\lnot}(\Undef) = \Undef$. وإذا كان أحد مقترني الاقتران غير
محدد والآخر صادقًا، لم يتعين الاقتران: فقد ينتهي صادقًا أو كاذبًا
بحسب ما ينتهي إليه المقترن المستقبلي. ولذلك
\[
    \tf{\land}(\True, \Undef) =
\tf{\land}(\Undef, \True) =
\Undef.
\]
وأما إن كان المقترن الآخر كاذبًا، امتنع أن يصير الاقتران صادقًا،
فيكون
\[\tf{\land}(\False, \Undef) =
\tf{\land}(\Undef, \False) = \False.\]
وإذا كانت الوسائط قيمًا متعينة، أي~$\True$ أو~$\False$، بقيت سائر
القيم كما هي في المنطق الكلاسيكي.

وأما الشرط ففيه تفصيل أكثر. لتكن~$q$ قضية مستقبلية عرضية. إذا كان
$p$ كاذبًا، صدق~$p \lif q$ كيفما انتهى~$q$؛ فنضع
$\tf{\lif}(\False, \Undef) = \True$. وإذا كان~$p$ صادقًا، صدق
$q \lif p$ كيفما انتهى~$q$؛ ولذلك
$\tf{\lif}(\Undef, \True) = \True$. وإذا كان~$p$ صادقًا، جاز أن
ينتهي~$p \lif q$ إلى الصدق أو الكذب، فنضع
$\tf{\lif}(\True, \Undef) = \Undef$. وعلى نظيره، إذا كان~$p$
كاذبًا، جاز أن ينتهي~$q \lif p$ إلى الوجهين؛ ومن ثم
$\tf{\lif}(\Undef, \False) = \Undef$. بقي أن يكون~$p$ و$q$
كلاهما مستقبليًا عرضيًا. ومقتضى الدافع الأول أن تكون القيمة حينئذ
$\Undef$؛ إلا أن هذا يسلب~$!A \lif !A$ كونه تحصيلًا منطقيًا.
وقد رضي~\L ukasiewicz بالتخلي عن~$!A \lor \lnot !A$ وعن
$\lnot(!A \land \lnot !A)$، لكنه لم يرض بسهولة بالتخلي عن
$!A \lif !A$؛ فاشترط~$\tf{\lif}(\Undef, \Undef) = \True$.

\begin{defn}\ollabel{def:lukasiewicz}
مصفوفة منطق~\L ukasiewicz الثلاثي القيم هي:
\begin{enumerate}
  \item الجزء الخالي من الثوابت~$\Lang L_0^{-}$ من لغة القضايا القياسية، ذو الروابط
  $\lnot$ و$\land$ و$\lor$ و$\lif$.
  % تصحيح مصدر مشترك (MVL-L0-I1): حُددت اللغة الخالية من lfalse التي تفسر المصفوفة جميع روابطها.
  \item مجموعة قيم الصدق~$V = \{\True, \Undef, \False\}$.
  \item القيمة المميزة وحدها هي~$\True$؛ أي~$V^+ = \{\True\}$.
  \item دوال الصدق مبينة في الجداول الآتية:
  \begin{center}
    \begin{tabular}{c|c}
      $\tf{\lnot}$ & \\
      \hline
      $\True$ & $\False$ \\
      $\Undef$ & $\Undef$ \\
      $\False$ & $\True$
    \end{tabular}
    \quad
    \begin{tabular}{c|ccc}
      $\tf{\land}[\LogLuk[3]]$ & $\True$ & $\Undef$ & $\False$ \\
      \hline
      $\True$ & $\True$ & $\Undef$ & $\False$ \\
      $\Undef$ & $\Undef$ & $\Undef$ & $\False$\\
      $\False$ & $\False$ & $\False$ & $\False$
    \end{tabular}
    \\[2ex]
    \begin{tabular}{c|ccc}
      $\tf{\lor}[\LogLuk[3]]$ & $\True$ & $\Undef$ & $\False$ \\
      \hline
      $\True$ & $\True$ & $\True$ & $\True$ \\
      $\Undef$ & $\True$ & $\Undef$ & $\Undef$ \\
      $\False$ & $\True$ & $\Undef$ & $\False$
    \end{tabular}
    \quad
    \begin{tabular}{c|ccc}
      $\tf{\lif}[\LogLuk[3]]$ & $\True$ & $\Undef$ & $\False$ \\
      \hline
      $\True$ & $\True$ & $\Undef$ & $\False$ \\
      $\Undef$ & $\True$ & $\True$ & $\Undef$  \\
      $\False$ & $\True$ & $\True$ & $\True$
    \end{tabular}
  \end{center}
\end{enumerate}
\end{defn}

ومن السهل أن كل~!!{formula}~$!A$ لا يدخل في تركيبها إلا~$\lnot$
و$\land$ و$\lor$ تأخذ~$\Undef$ إذا أسند~$\Undef$ إلى جميع
!!{propositional variable}s فيها. فلا يكون، مثلًا، التحصيلان
الكلاسيكيان~$p \lor \lnot p$ و$\lnot(p \land \lnot p)$ تحصيلين
في~$\LogLuk[3]$؛ إذ~$\pValue{v}(!A) = \Undef$ متى كان
$\pAssign v(p) = \Undef$.

وأما~!!{valuation}s التي يكون فيها~$\pAssign v(p) = \True$ أو
$\False$، فتوافق فيها~$\pValue v(!A)$ قيمة الصدق الكلاسيكية.

\begin{prop}
  إذا كان~$\pAssign v(p) \in \{\True, \False\}$ لكل~$p$ في~$!A$،
  كان~$\pValue v[\LogLuk[3]](!A) = \pValue v[\LogCL](!A)$.
\end{prop}

\begin{prob}\label{mvl:thr:luk:prob:luk-iff} إذا عرّفنا
  $\pValue v(!A \liff !B) = \pValue v((!A \lif !B) \land (!B \lif
  !A))$ في~$\LogLuk[3]$، فما جدول صدق~$\liff$؟
\end{prob}

وكثير من التحصيلات الكلاسيكية تحصيلات أيضًا في~\LogLuk[3]، ومنها
$\lnot p \lif (p \lif q)$. ويمكن التثبت منها، كما في المنطق
الكلاسيكي، بجدول الصدق:
\begin{center}
\begin{tabular}{cc|cccccc}
  $p$ & $q$ & $\lnot$ & $p$ & $\lif$ & $\smash{(}p$ & $\lif$ & $q\smash{)}$ \\ \hline
  \True & \True & \False & \True & \True & \True & \True & \True \\
  \True & \Undef & \False & \True & \True & \True & \Undef & \Undef \\
  \True & \False & \False & \True & \True & \True & \False & \False \\
  \Undef & \True & \Undef & \Undef & \True & \Undef & \True & \True \\
  \Undef & \Undef & \Undef & \Undef & \True & \Undef & \True & \Undef \\
  \Undef & \False & \Undef & \Undef & \True & \Undef & \Undef & \False \\
  \False & \True & \True & \False & \True & \False & \True & \True \\
  \False & \Undef & \True & \False & \True & \False & \True & \Undef \\
  \False & \False & \True & \False & \True & \False & \True & \False \\
\end{tabular}
\end{center}

\begin{prob}
  بيّن أن الصيغ الآتية تحصيلات منطقية في~\LogLuk[3]:
  \begin{enumerate}
    \item $p \lif (q \lif p)$
    \item\label{mvl:thr:luk:prob:luk-taut-2} $\lnot(p \land q) \liff (\lnot p \lor \lnot q)$
    \item\label{mvl:thr:luk:prob:luk-taut-3} $\lnot(p \lor q) \liff (\lnot p \land \lnot q)$
  \end{enumerate}
  (في \olref[mvl][thr][luk]{prob:luk-taut-2} و
  \olref[mvl][thr][luk]{prob:luk-taut-3}، خذ~$!A \liff !B$
  اختصارًا لـ~$(!A \lif !B) \land (!B \lif !A)$، أو استعمل جوابك
  عن \olref[mvl][thr][luk]{prob:luk-iff}.)
\end{prob}

\begin{prob}
  بيّن أن التحصيلات الكلاسيكية الآتية لا تكون تحصيلات في~\LogLuk[3]:
  \begin{enumerate}
    \item $(\lnot p \land p) \lif q$
    \item $((p \lif q) \lif p) \lif p$
    \item $(p \lif (p \lif q)) \lif (p \lif q)$
  \end{enumerate}
\end{prob}

وقد يحمل ذلك على الظن أن كل تحصيل كلاسيكي، وإن لم يكن تحصيلًا في
$\LogLuk[3]$، لا بد أن يأخذ في كل~!!{valuation} واحدة من القيمتين
$\True$ و$\Undef$. وليس الأمر كذلك، وهذه الصيغة مثال مضاد:
\[
  \lnot(p \lif \lnot p) \lor \lnot(\lnot p \lif p)
\]
فإنها تأخذ~$\False$ إذا أخذ~$p$ القيمة~$\Undef$.

\begin{prob}
  أي العلاقات الآتية تصح في منطق~\L ukasiewicz؟ ضع لكل واحدة جدول صدق.
  \begin{enumerate}
    \item $p, p \lif q \Entails q$
    \item $\lnot\lnot p \Entails p$
    \item $p \land q \Entails p$
    \item $p \Entails p \land p$
    \item $p \Entails p \lor q$
  \end{enumerate}
\end{prob}

وكان رجاء~\L ukasiewicz أن يبني على نسقه الثلاثي منطقًا للإمكان،
بأن يزيد الرابط الأحادي~$\Diamond !A$، أي «$!A$ ممكنة»، ونظيره
$\Box !A$، أي «$!A$ ضرورية»:
\begin{center}
  \begin{tabular}{c|c}
    $\tf{\Diamond}$ & \\
    \hline
    $\True$ & $\True$ \\
    $\Undef$ & $\True$ \\
    $\False$ & $\False$
  \end{tabular}
  \quad
  \begin{tabular}{c|c}
    $\tf{\Box}$ & \\
    \hline
    $\True$ & $\True$ \\
    $\Undef$ & $\False$ \\
    $\False$ & $\False$
  \end{tabular}
\end{center}
فمؤدى الجدولين أن~$p$ ممكن إذا وفقط إذا لم يتعين كاذبًا بالفعل،
وأن~$p$ ضروري إذا وفقط إذا تعين صادقًا بالفعل.

\begin{prob}
  بيّن أن~$\Box p \liff \lnot\Diamond \lnot p$ و
  $\Diamond p \liff \lnot \Box \lnot p$ تحصيلان منطقيان في
  $\LogLuk[3]$ بعد توسيعه بجدولي صدق~$\Box$ و$\Diamond$.
\end{prob}

لكن قصور هذا المنطق الجهي المقترح يظهر سريعًا. فمهما آل الأمر، يمتنع
أن يتبين~$p \land \lnot p$ صادقًا؛ فهو مستحيل، وإن لم يتعين الآن
فكان غير محدد. فينبغي إذن أن تكون~$\lnot \Diamond(p \land \lnot p)$
تحصيلًا منطقيًا. غير أن~$\pAssign v(p) = \Undef$ يجعل
% تصحيح مصدر (0394-I2): يعطي جدولا النفي والإمكان هنا القيمة الكاذبة.
$\pValue v(\lnot \Diamond(p \land \lnot p)) = \False$.
فقد أصاب~\L ukasiewicz في أن قيمتي صدق لا تفيان بتمييزات الجهة،
كالإمكان والضرورة، لكن مجرد إضافة قيمة ثالثة لا يفي بها أيضًا.

\end{document}
